\clearpage\section{MOS transistor schematic}
These editable vector schematics connect equal node names within a circuit. Primitive-cell names are local pins; the hierarchy defines their connections. B denotes the body terminal; 0 is global ground. Each panel identifies its netlist and scope; measurement sources are connected as described.
All ten MOS devices of \texttt{05\_complete\_reference.sp} are shown in three connected panels. All device lengths are $0.3\,\mu$m. Widths below are in $\mu$m.
\par
\begin{figure}[!ht]
\centering
\begin{circuitikz}[font=\small]
\node[pmos,bulk] (M1) at (0,3) {};
\node[anchor=west,font=\scriptsize] at ($(M1.center)+(.55,.40)$) {M1};
\draw (M1.bulk) -- ++(.30,0) node[right,font=\scriptsize] {B=vdd};
\draw (M1.G) -- ++(-.45,0) node[left,font=\small] {\texttt{bp}};
\node[nmos,bulk] (M3) at (0,0.7) {};
\node[anchor=west,font=\scriptsize] at ($(M3.center)+(.55,.40)$) {M3};
\draw (M3.bulk) -- ++(.30,0) node[right,font=\scriptsize] {B=0};
\draw (M3.G) -- ++(-.45,0) node[left,font=\small] {\texttt{nx}};
\node[pmos,bulk] (M2) at (4.5,3) {};
\node[anchor=west,font=\scriptsize] at ($(M2.center)+(.55,.40)$) {M2};
\draw (M2.bulk) -- ++(.30,0) node[right,font=\scriptsize] {B=vdd};
\draw (M2.G) -- ++(-.45,0) node[left,font=\small] {\texttt{bp}};
\node[nmos,bulk] (M4) at (4.5,0.7) {};
\node[anchor=west,font=\scriptsize] at ($(M4.center)+(.55,.40)$) {M4};
\draw (M4.bulk) -- ++(.30,0) node[right,font=\scriptsize] {B=0};
\draw (M4.G) -- ++(-.45,0) node[left,font=\small] {\texttt{nx}};
\draw (M1.S) -- ++(0,0.5) node[above,font=\small] {\texttt{rs}};
\draw (M1.D) -- (M3.D) node[midway,above right,font=\small] {\texttt{nx}};
\draw (M3.S) -- ++(0,-0.4) node[below,font=\small] {\texttt{0}};
\draw (M2.S) -- ++(0,0.5) node[above,font=\small] {\texttt{vdd}};
\draw (M2.D) -- (M4.D) node[midway,above right,font=\small] {\texttt{bp}};
\draw (M4.S) -- ++(0,-0.4) node[below,font=\small] {\texttt{0}};
\draw (8.2,4) node[above] {\texttt{vdd}} to[R,l={RB}] ++(0,-1.7) node[below] {\texttt{rs}};
\draw (8.2,1) node[above] {\texttt{nx}} to[C,l={CNX}] ++(0,-1.7) node[below] {\texttt{0}};
\draw (11.0,1) node[above] {\texttt{bp}} to[C,l={CBP}] ++(0,-1.7) node[below] {\texttt{0}};
\end{circuitikz}
\caption{Self-biased reference core: M1=115, M2=22, M3=M4=4.85; RB=rbias. CNX=CBP=20 fF.}
\end{figure}
\par
\par
\begin{figure}[!ht]
\centering
\begin{circuitikz}[font=\small]
\node[pmos,bulk] (MS1) at (0,3) {};
\node[anchor=west,font=\scriptsize] at ($(MS1.center)+(.55,.40)$) {MS1};
\draw (MS1.bulk) -- ++(.30,0) node[right,font=\scriptsize] {B=vdd};
\draw (MS1.G) -- ++(-.45,0) node[left,font=\small] {\texttt{bp}};
\node[nmos,bulk] (MS2) at (0,0.7) {};
\node[anchor=west,font=\scriptsize] at ($(MS2.center)+(.55,.40)$) {MS2};
\draw (MS2.bulk) -- ++(.30,0) node[right,font=\scriptsize] {B=0};
\draw (MS2.G) -- ++(-.45,0) node[left,font=\small] {\texttt{bp}};
\node[pmos,bulk] (MS3) at (4.5,3) {};
\node[anchor=west,font=\scriptsize] at ($(MS3.center)+(.55,.40)$) {MS3};
\draw (MS3.bulk) -- ++(.30,0) node[right,font=\scriptsize] {B=vdd};
\draw (MS3.G) -- ++(-.45,0) node[left,font=\small] {\texttt{st}};
\draw (MS1.S) -- ++(0,0.5) node[above,font=\small] {\texttt{vdd}};
\draw (MS1.D) -- (MS2.D) node[midway,above right,font=\small] {\texttt{st}};
\draw (MS2.S) -- ++(0,-0.4) node[below,font=\small] {\texttt{0}};
\draw (MS3.S) -- ++(0,0.5) node[above,font=\small] {\texttt{vdd}};
\draw (MS3.D) -- ++(0,-0.4) node[below,font=\small] {\texttt{inj}};
\draw (8.2,3.8) node[above] {\texttt{inj}} to[R,l={RSTART}] ++(0,-1.7) node[below] {\texttt{nx}};
\draw (8.2,0.8) node[above] {\texttt{st}} to[C,l={CST}] ++(0,-1.7) node[below] {\texttt{0}};
\end{circuitikz}
\caption{Startup network: MS1=2, MS2=0.2, MS3=10; CST=20 fF, RSTART=rstart.}
\end{figure}
\par
\clearpage
\par
\begin{figure}[!ht]
\centering
\begin{circuitikz}[font=\small]
\node[pmos,bulk] (M5) at (0,5) {};
\node[anchor=west,font=\scriptsize] at ($(M5.center)+(.55,.40)$) {M5};
\draw (M5.bulk) -- ++(.30,0) node[right,font=\scriptsize] {B=vdd};
\draw (M5.G) -- ++(-.45,0) node[left,font=\small] {\texttt{bp}};
\node[pmos,bulk] (MP) at (0,2.8) {};
\node[anchor=west,font=\scriptsize] at ($(MP.center)+(.55,.40)$) {MP};
\draw (MP.bulk) -- ++(.30,0) node[right,font=\scriptsize] {B=vdd};
\draw (MP.G) -- ++(-.45,0) node[left,font=\small] {\texttt{vref}};
\node[nmos,bulk] (MN) at (0,0.6) {};
\node[anchor=west,font=\scriptsize] at ($(MN.center)+(.55,.40)$) {MN};
\draw (MN.bulk) -- ++(.30,0) node[right,font=\scriptsize] {B=0};
\draw (MN.G) -- ++(-.45,0) node[left,font=\small] {\texttt{div}};
\draw (M5.S) -- ++(0,0.5) node[above,font=\small] {\texttt{vdd}};
\draw (M5.D) -- (MP.S) node[midway,above right,font=\small] {\texttt{top}};
\draw (MP.D) -- (MN.D) node[midway,above right,font=\small] {\texttt{vref}};
\draw (MN.S) -- ++(0,-0.4) node[below,font=\small] {\texttt{0}};
\draw (4,5.6) node[above]{\texttt{top}} to[R,l=R3] (4,3.2) node[circ,label=right:\texttt{div}]{} to[R,l=R4] (4,.4) node[below]{0};
\draw (8,3.5) node[above] {\texttt{vref}} to[C,l={CL}] ++(0,-1.7) node[below] {\texttt{0}};
\end{circuitikz}
\caption{Reference output: M5=MP=22, MN=4.85; R3=rupper, R4=1 k$\Omega$, CL=1 pF.}
\end{figure}
\par
M1's source is rs and MP's source is top; their bodies remain at vdd. These are separate body/source terminals, not body-to-source shorts. Baseline rbias=1 k$\Omega$, rstart=1 m$\Omega$, rupper=5 k$\Omega$, vdd=1 V. The MOS-only divider and startup experiments use their own source decks; these panels show the combined reference.
